% 本文件由 assemble.py 按 _work/plans/ch16.json 逐字组装，请勿手改（改 plan 或 blocks 后重新生成）。
\chapter{热学}
\label{ch:16}

\section{气体状态变化图像（$V$-$T$、$p$-$V$）}
% ---- block 185-03 (原稿第185页) kind=problem ----
% 粘贴的印刷题（紫色纸），右端被页边截断；只有题干，无解答
\begin{problem}[粘贴印刷题 24(1)]
24．(1)如图所示，一定质量的理想气体由状态 $a$ 开始，经历 $ab$、$bc$、$cd$、$da$ 四个过程回到原状态，其 $V-T$……% 此行右端被页边截断
刚好是一个圆，气体在 $a$、$b$、$c$、$d$ 四个状态的体积分别为 $V_a=V_c$、$V_b=V_2$、$V_d=V_1$，温度分别为 $T_a=T_1$、$T_b=$……% 此行右端被页边截断
$T_c=T_2$，压强分别为 $p_a$、$p_b$、$p_c$、$p_d$，其中 $V_2=5V_1$，$T_2=5T_1$。下列说法正确的是（\qquad）

\medskip
\noindent\textbf{手写旁注（右下，下半被粘贴纸边缘遮挡）：}\unclear{$a=3V$}；\unclear{$\dfrac{10}{\ldots}$}，\unclear{.2}% CHECK: 被遮挡，可能是 V_a=3V_1 或 9=3V（与上方 9=2V 同类）
\end{problem}

% ---- block 224-03 (原稿第224页) kind=problem ----
\begin{problem}
\notefig[0.25]{figures/p224_a.png}
\notefig[0.25]{figures/p224_b.png}
\begin{solution}
$1\unit{atm}=1\times10^5\unit{Pa}$

$W=$\unclear{…}$(S_1-S_2)\times(-1)$ % 原稿：“=”后有一处被涂划的字迹

$W=-4000\unit{J}$\quad $\Delta U=Q+W=0$\quad 故吸热 $4000\unit{J}$
\end{solution}
\end{problem}
功的正负\quad 在 顺吸 逆放

